\chapter{What was predicted, and what was wrong}
\label{ch:predictions}

\section{Why the predictions are here at all}

A prediction committed after the measurement is not a prediction. Every
prediction in this project was written into a file, committed, and a continuous
integration job verifies that the prediction commit is an ancestor of the commit
carrying the result files it predicts about. That check is the only mechanism
that makes the fourth law more than an intention.

The predictions are printed beside the outcomes because that is the only form in
which a prediction is worth anything. A report that showed only the outcomes
would be indistinguishable from a report whose author had decided afterwards
which results were interesting.

\begin{table}[htbp]
  \centering
  \caption{The eight claims registered before the benchmark ran, the number that
  settled each one, and the outcome. The measured column is read from the result
  files by the table generator, so a prediction cannot be scored generously by
  whoever writes the report.}
  \label{tab:predictions}
  \begin{longtable}{p{.34\textwidth}p{.30\textwidth}p{.28\textwidth}}
\toprule
\textbf{Registered before the run} & \textbf{Measured} & \textbf{Outcome} \\
\midrule
\endfirsthead
\toprule
\textbf{Registered before the run} & \textbf{Measured} & \textbf{Outcome} \\
\midrule
\endhead
\bottomrule
\endfoot
The language model beats the n-gram model on the hostile tier & llm 0.5265 against ngram 0.5409 & wrong: it lost \\  % results: experiments/results/bench/2026-08-26T20-32-48Z_bench_6457ac7/benchmark.json experiments/results/test/2026-08-26T20-32-48Z_p6-llm_6457ac7/metrics.json
The language model beats the n-gram model on the held out locale & llm 0.6045 against ngram 0.5637 & held, and not certified after correction \\  % results: experiments/results/bench/2026-08-26T20-32-48Z_bench_6457ac7/benchmark.json experiments/results/test/2026-08-26T20-32-48Z_p6-llm_6457ac7/metrics.json
The language model does not beat the rule table on the held out locale & llm 0.6045 against rules 0.7276 & held \\  % results: experiments/results/bench/2026-08-26T20-32-48Z_bench_6457ac7/benchmark.json experiments/results/test/2026-08-26T20-32-48Z_p6-llm_6457ac7/metrics.json
Latency differs by three orders of magnitude or more & p95 674727.4 against 134.0 & held \\  % results: experiments/results/bench/2026-08-26T20-32-48Z_bench_6457ac7/benchmark.json experiments/results/test/2026-08-26T20-32-48Z_p6-llm_6457ac7/metrics.json
Schema failures above zero and below one request in twenty & 0 retried of 480 requests & wrong: exactly none \\  % results: experiments/results/bench/2026-08-26T20-32-48Z_bench_6457ac7/benchmark.json experiments/results/test/2026-08-26T20-32-48Z_p6-llm_6457ac7/metrics.json
The n-gram model stays the right default & ngram macro-F1 0.7394 against llm 0.6353 & held, for a stronger reason than the one registered \\  % results: experiments/results/bench/2026-08-26T20-32-48Z_bench_6457ac7/benchmark.json experiments/results/test/2026-08-26T20-32-48Z_p6-llm_6457ac7/metrics.json
tel-national read as tel occurs; username read as email does not & tel-national as tel 14; username as email 0 & held in both directions \\  % results: experiments/results/bench/2026-08-26T20-32-48Z_bench_6457ac7/benchmark.json experiments/results/test/2026-08-26T20-32-48Z_p6-llm_6457ac7/metrics.json
Abstention lands between the two classical engines & llm 0.0255, ngram 0.0363, rules 0.2620 & wrong: below both \\  % results: experiments/results/bench/2026-08-26T20-32-48Z_bench_6457ac7/benchmark.json experiments/results/test/2026-08-26T20-32-48Z_p6-llm_6457ac7/metrics.json
\end{longtable}
\par\smallskip\noindent{\footnotesize Source: \rpath{experiments/results/bench/2026-08-26T20-32-48Z_bench_6457ac7/benchmark.json}}
\par\noindent{\footnotesize Source: \rpath{experiments/results/test/2026-08-26T20-32-48Z_p6-llm_6457ac7/metrics.json}}

\end{table}

\section{Four held and four did not}

\textbf{The four that failed are the more interesting half}, and two of them
should be read first.

\textbf{Schema compliance was predicted to fail sometimes, and never failed.}
The registered claim was that the retry rate would be above zero and below one
request in twenty. It was exactly zero across every request of the run. The
mechanism appears to be that passing the response schema with the request lets
the server constrain decoding, which removes the failure mode rather than
reducing it. The retry ladder was built for a problem that, on this serving
route, does not arise. It is still there and still tested, because the next
serving route may not have that property.

\textbf{Abstention went the wrong way.} It was registered to land between the two
classical engines, on the reasoning that a model instructed to decline would
decline more than a model with no such instruction and less than a rule table
whose default branch abstains by construction. It landed below both. Instructing
a model to decline is not equivalent to giving it a threshold, and this is the
measurement of that difference.

\textbf{The hostile tier prediction was wrong in the direction that matters
most.} The hostile tier is where labels are stripped and only obfuscated
identifiers survive, which is where broad pretrained knowledge was supposed to
pay for itself. The language model lost there too.

\textbf{The confusion prediction was wrong in direction as well as in detail.}
The failure was expected to be over claiming: fields that are not personal data
read as though they were. The largest single confusion is the reverse, fields
whose answer key is \code{UNKNOWN} read as the label for a control with no
autofill meaning. The model declines more readily than predicted on the fields
where declining is wrong, while declining less readily overall than either
classical engine. Table~\ref{tab:predicted-pairs} in Appendix~\ref{app:tables}
gives the counts for every confusion pair the specification named in advance, and
Table~\ref{tab:confusions} gives the largest confusions each engine actually
makes.

\section{The pattern across every phase of this project}

Read across all four prediction files rather than only this one, a pattern is
visible and it is not flattering to the intuition that produced the predictions.

\textbf{Every prediction that failed was about a model. Every prediction that
held was about the system.} The predictions that held were the ones where the
mechanism was already understood: page load dominates because a browser is slow;
the n-gram engine is slower per field than the rule table because it enumerates
n-grams in interpreted code; latency differs by orders of magnitude because one
engine is a regular expression match and the other is a forward pass. The ones
that failed were the ones where a mechanism was assumed rather than measured: the
model will generalise better across locales; broad knowledge will help on
degraded markup; instructing a model to abstain will make it abstain.

Two earlier phases are worth listing beside this one, because the pattern is what
makes them worth listing.

At the phase that built the n-gram model, the prediction that the derived
decision threshold would sit \emph{below} the rule engine's confident band was
wrong in the other direction: it sits well above it. That single wrong prediction
is the reason the two engines carry separate threshold blocks, and it is the most
consequential wrong prediction in the project.

At the first evaluation phase, three predictions about the model winning all
failed, and a fourth about which confusion pairs would appear failed completely:
not one of the four pairs named in advance occurred even once, in either
direction, for either engine.

\section{The temptation that was registered against}

The prediction file for this benchmark records, in advance, the specific
temptation that a three engine comparison creates: to correct across the one
engine pair the project cared about rather than across all three. It also records
the consequence of refusing that, which is that every comparison between the two
classical engines would come out with a larger adjusted p than before.

That is the shape a useful pre-registration takes. It is not a forecast of the
result. It is a written commitment to the analysis choices that would otherwise
be made after seeing it.
